\documentclass[12pt]{article}
\usepackage[utf8]{inputenc}
\usepackage[T1]{fontenc}
\usepackage{hyperref}
\title{Tuxin\'awa Language Resources for K-12 Curriculum Development}
\author{Compiled by Claude Sonnet 5 (claude-sonnet-5) for Andr\'as Kornai}
\date{}
\begin{document}
\maketitle

\section{Introduction}
Tuxin\'awa (ISO 639-3 tux; Glottocode tuxi1238) is an extinct Panoan language of Acre, Brazil. It is reported as extinct by the 1970s and resembled Yaminawa dialects. Search results turned up no native-language teaching materials. The name Tuxinawa is also used as an alternate name for the living Cashinahua (Kaxinaw\'a) language, which is a different language and must not be confused with this one.

\section{Primary Educational Resources}
None were found. No primers, textbooks or school programmes for Tuxin\'awa appeared in search results.

\section{Linguistic and Reference Materials}
\begin{itemize}
\item Wikipedia entry on the Tuxinawa language (classification and status): \url{https://en.wikipedia.org/wiki/Tuxinawa_language}
\item ASJP database entry with a short wordlist. The listing indicates the data come from a 1931 source by Braulino de Carvalho: \url{https://asjp.clld.org/languages/TUXINAWA}
\item Endangered Languages Project entry: \url{https://www.endangeredlanguages.com/lang/1846?hl=es}
\item Overview of the Panoan family, useful for comparative work: \url{https://en.wikipedia.org/wiki/Panoan_languages}
\end{itemize}

\section{Religious and Community Resources}
None were found.

\section{Digital and Audio Resources}
None were found.

\section{Curriculum Development Context}
The language is extinct and documentation is limited to a short wordlist. A K-12 heritage curriculum could not be based on Tuxin\'awa itself. Any work would have to treat the language as a historical topic and use closely related living Panoan languages for comparison. Because the name Tuxinawa is shared with Cashinahua, any Cashinahua materials should be confirmed as a separate language before use. This dossier is short because little exists.
\end{document}